\documentclass[11pt,a4paper]{article}
\usepackage[utf8]{inputenc}
\usepackage[vietnamese]{babel}
\usepackage{amsmath,amssymb,amsfonts}
\usepackage{geometry}
\geometry{a4paper, top=1.2cm, bottom=1.2cm, left=1.4cm, right=1.4cm}
\usepackage{tikz}
\usetikzlibrary{calc, intersections, angles, arrows.meta, shapes.geometric, positioning}
\usepackage{tcolorbox}
\tcbuselibrary{skins, breakable}
\usepackage{tabularx}
\usepackage{enumitem}
\usepackage{fancyhdr}

\pagestyle{fancy}
\fancyhf{}
\setlength{\headheight}{15pt}
\lhead{\small\textbf{TOÁN 9+ | BẺ KHÓA HÌNH HỌC VÀO LỚP 10}}
\rhead{\small\textbf{BÀI 02: BỔ ĐỀ TRỰC TÂM \& CỰC TRỊ}}
\cfoot{\thepage}

% Custom tcolorbox styles
\tcbset{
    boxrule=0.8pt,
    arc=2.5mm,
    fonttitle=\bfseries\sffamily,
    coltitle=white
}

\newtcolorbox{modelbox}[1]{
    colback=blue!4!white,
    colframe=blue!75!black,
    title=#1
}

\newtcolorbox{formulabox}[1]{
    colback=orange!3!white,
    colframe=orange!85!black,
    title=#1
}

\newtcolorbox{pedagogybox}[1]{
    colback=teal!4!white,
    colframe=teal!80!black,
    title=#1
}

\newtcolorbox{solutionbox}[1]{
    colback=green!4!white,
    colframe=green!60!black,
    title=#1
}

\newtcolorbox{exercisebox}[1]{
    colback=purple!4!white,
    colframe=purple!75!black,
    title=#1
}

\begin{document}

% ==============================================================
% TRANG 1: ĐỀ BÀI VÀ HÌNH VẼ TIKZ TRỰC QUAN
% ==============================================================
\begin{center}
    {\Large\bfseries\color{blue!80!black} TUYỆT PHẨM HÌNH HỌC 9+ VÀO LỚP 10 CHUYÊN \& CÔNG LẬP TOP ĐẦU}\\[2mm]
    {\large\bfseries BÀI HÌNH 02: BỔ ĐỀ TRỰC TÂM KINH ĐIỂN \& PHƯƠNG PHÁP CỰC TRỊ CAUCHY}\\[1mm]
    \textit{(Trích xuất và phát triển từ Đề 02 - Tuyển tập 55 Đề Thi Tuyển Sinh Vào Lớp 10)}
\end{center}

\vspace{1mm}

\begin{modelbox}{I. ĐỀ BÀI TOÁN GỐC (TRÍCH ĐỀ 02)}
Cho đường tròn $(O)$, bán kính $R (R>0)$ và dây cung $BC$ cố định. Một điểm $A$ chuyển động trên cung lớn $BC$ sao cho tam giác $ABC$ có ba góc nhọn. Kẻ các đường cao $AD, BE$ của tam giác $ABC$ cắt nhau tại $H$.
\begin{enumerate}[label=\textbf{\alph*)}]
    \item Chứng minh rằng tứ giác $DHEC$ nội tiếp một đường tròn.
    \item Kẻ đường kính $AM$ của đường tròn $(O)$ và $OI$ vuông góc với $BC$ tại $I$. Chứng minh rằng $I$ là trung điểm của $HM$.
    \item \textbf{(Phân hóa 9+):} Khi $BC$ cố định, xác định vị trí của $A$ trên đường tròn $(O)$ để tích $DH \cdot DA$ đạt giá trị lớn nhất.
\end{enumerate}
\end{modelbox}

\vspace{1mm}

\begin{center}
\begin{tikzpicture}[line join = round, line cap=round, >=stealth, font=\footnotesize, scale=1.3]
    \def\r{3.0}
    \coordinate (O) at (0,0);
    \coordinate (A) at (105:\r);
    \coordinate (B) at (-155:\r);
    \coordinate (C) at (-25:\r);
    
    % Heights
    \coordinate (D) at ($(B)!(A)!(C)$);
    \coordinate (E) at ($(A)!(B)!(C)$);
    
    \path[name path=lineAD] (A) -- ($(A)!1.8!(D)$);
    \path[name path=lineBE] (B) -- ($(B)!1.8!(E)$);
    \path[name intersections={of=lineAD and lineBE, by=H}];
    
    % Diameter AM
    \coordinate (M) at ($(O)!-1!(A)$);
    
    % Midpoint I of BC
    \coordinate (I) at ($(B)!0.5!(C)$);
    
    % Intersection K of AD with (O)
    \coordinate (K) at ($2*(D)-(H)$);
    
    % Draw circle and triangle
    \draw[thick, blue!80!black] (O) circle (\r);
    \draw[thick] (A) -- (B) -- (C) -- cycle;
    
    % Heights
    \draw[blue!70!black] (A) -- (D);
    \draw[blue!70!black] (B) -- (E);
    \draw[red!80!black, thick] (D) -- (K);
    \draw[teal!80!black] (B) -- (K);
    
    % Diameter AM and parallelogram BHCM
    \draw[dashed, purple!80!black] (A) -- (M);
    \draw[teal!80!black, thick] (B) -- (M) -- (C);
    \draw[purple!80!black] (B) -- (H) -- (C);
    \draw[dashed, gray] (H) -- (M);
    \draw[dashed, gray] (O) -- (I);
    
    % Points
    \fill (O) circle (1.5pt) node[above right] {$O$};
    \fill (A) circle (1.5pt) node[above] {$A$};
    \fill (B) circle (1.5pt) node[below left] {$B$};
    \fill (C) circle (1.5pt) node[below right] {$C$};
    \fill (D) circle (1.5pt) node[below right] {$D$};
    \fill (E) circle (1.5pt) node[above right] {$E$};
    \fill (H) circle (1.5pt) node[left, red!85!black] {$H$};
    \fill (M) circle (1.5pt) node[below, purple!85!black] {$M$};
    \fill (I) circle (1.5pt) node[below, teal!85!black] {$I$};
    \fill (K) circle (1.5pt) node[below, red!85!black] {$K$};
    
    % Right angles
    \pic[draw, angle radius=2.5mm]{right angle = A--D--B};
    \pic[draw, angle radius=2.5mm]{right angle = B--E--C};
    \pic[draw, angle radius=2.5mm]{right angle = O--I--C};
    \pic[draw, purple!80!black, angle radius=2.5mm]{right angle = A--B--M};
    \pic[draw, purple!80!black, angle radius=2.5mm]{right angle = A--C--M};
\end{tikzpicture}
\end{center}

\newpage

% ==============================================================
% TRANG 2: SƠ ĐỒ CÔNG THỨC SUY LUẬN NGƯỢC THUẦN TÚY
% ==============================================================
\begin{formulabox}{II. SƠ ĐỒ CÔNG THỨC SUY LUẬN NGƯỢC THUẦN TÚY (LOGIC FORMULA TREE)}

\textbf{\color{orange!90!black} 1. CÂY SUY LUẬN CÔNG THỨC CHO CÂU B (HÌNH BÌNH HÀNH $BHCM$):}
\vspace{1mm}

\noindent
\begin{tikzpicture}[
    every node/.style={align=center, font=\footnotesize},
    box/.style={rectangle, draw=blue!75!black, fill=blue!5, thick, rounded corners=1.5mm, inner sep=4pt, text width=16.4cm},
    headbox/.style={rectangle, draw=red!80!black, fill=red!8, very thick, rounded corners=1.5mm, inner sep=4pt, text width=16.4cm},
    basebox/.style={rectangle, draw=teal!85!black, fill=teal!10, very thick, rounded corners=1.5mm, inner sep=4pt, text width=16.4cm},
    arr/.style={-{Stealth[length=2.5mm]}, very thick, red!80!black}
]

\node[headbox] (b1_1) {
    \textbf{\color{red!80!black} [ĐÍCH ĐẾN CÂU B]:} \quad $\boldsymbol{I \text{ LÀ TRUNG ĐIỂM CỦA } HM}$
};

\node[box, below=0.22cm of b1_1] (b1_2) {
    $\boldsymbol{\Uparrow} \quad \boldsymbol{BHCM \text{ là hình bình hành}} \quad \text{và } I \text{ là trung điểm } BC \ (OI \perp BC)$
};

\node[basebox, below=0.22cm of b1_2] (b1_3) {
    $\boldsymbol{\Uparrow} \quad \begin{cases} BH \parallel MC \quad (\text{cùng } \perp AC \text{ do } BE \perp AC \text{ và } \widehat{ACM}=90^\circ) \\ CH \parallel MB \quad (\text{cùng } \perp AB \text{ do } CF \perp AB \text{ và } \widehat{ABM}=90^\circ) \end{cases} \quad \boldsymbol{\Longleftarrow} \quad \boldsymbol{AM \text{ là đường kính của } (O)}$
};

\draw[arr] (b1_3) -- (b1_2);
\draw[arr] (b1_2) -- (b1_1);

\end{tikzpicture}

\vspace{2.5mm}
\hrule
\vspace{2.5mm}

\textbf{\color{orange!90!black} 2. CÂY SUY LUẬN CÔNG THỨC CHO CÂU C PHÂN HÓA 9+ ($\max(DH \cdot DA)$):}
\vspace{1mm}

\noindent
\begin{tikzpicture}[
    every node/.style={align=center, font=\footnotesize},
    box/.style={rectangle, draw=blue!75!black, fill=blue!5, thick, rounded corners=1.5mm, inner sep=3.5pt, text width=16.4cm},
    headbox/.style={rectangle, draw=red!80!black, fill=red!10, very thick, rounded corners=1.5mm, inner sep=4pt, text width=16.4cm},
    actionbox/.style={rectangle, draw=orange!85!black, fill=orange!10, thick, rounded corners=1.5mm, inner sep=3.5pt, text width=16.4cm},
    basebox/.style={rectangle, draw=teal!85!black, fill=teal!12, very thick, rounded corners=1.5mm, inner sep=4pt, text width=16.4cm},
    arr/.style={-{Stealth[length=2.5mm]}, very thick, red!80!black}
]

\node[headbox] (b2_1) {
    \textbf{\color{red!80!black} [ĐÍCH ĐẾN CỰC TRỊ]:} \quad $\boldsymbol{\max(DH \cdot DA) = \dfrac{BC^2}{4}}$\\[0.5mm]
    $\boldsymbol{\iff A \text{ LÀ ĐIỂM CHÍNH GIỮA CỦA CUNG LỚN } BC}$
};

\node[box, below=0.22cm of b2_1] (b2_2) {
    $\boldsymbol{\Uparrow} \quad \text{Dấu “=” xảy ra} \iff \boldsymbol{DB = DC} \iff D \equiv I \iff \boldsymbol{\triangle ABC \text{ cân tại } A}$
};

\node[actionbox, below=0.22cm of b2_2] (b2_3) {
    $\boldsymbol{\Uparrow} \quad \text{\textbf{BẤT ĐẲNG THỨC CAUCHY:}}\quad \boldsymbol{DB \cdot DC \le \left(\dfrac{DB+DC}{2}\right)^2 = \left(\dfrac{BC}{2}\right)^2 = \dfrac{BC^2}{4}} \quad (BC \text{ cố định})$
};

\node[box, below=0.22cm of b2_3] (b2_4) {
    $\boldsymbol{\Uparrow} \quad \text{Quy đổi tích biến thiên về tích cố định:}\quad \boldsymbol{DH \cdot DA = DK \cdot DA = DB \cdot DC}$
};

\node[box, below=0.22cm of b2_4] (b2_5) {
    $\boldsymbol{\Uparrow} \quad \begin{cases} \boldsymbol{DK \cdot DA = DB \cdot DC} \quad (\triangle DBK \sim \triangle DAC \text{ g-g theo định lý phương tích}) \\ \boldsymbol{DH = DK} \quad (\triangle HBK \text{ cân tại } B \text{ có đường cao } BD \text{ là trung trực}) \end{cases}$
};

\node[box, below=0.22cm of b2_5] (b2_6) {
    $\boldsymbol{\Uparrow} \quad \boldsymbol{\widehat{HBD} = \widehat{KBD} = \widehat{KAC} = 90^\circ - \widehat{C}} \quad \left(\text{cùng chắn cung } \stackrel{\frown}{KC} \text{ của } (O) \text{ và cùng phụ } \widehat{C}\right)$
};

\node[basebox, below=0.22cm of b2_6] (b2_7) {
    $\boldsymbol{\Uparrow} \quad \text{\textbf{KÉO DÀI ĐƯỜNG CAO }} \boldsymbol{AD} \text{\textbf{ CẮT }} \boldsymbol{(O)} \text{\textbf{ TẠI ĐIỂM THỨ HAI }} \boldsymbol{K}$\\
    \vspace{0.5mm}
    \textbf{\color{teal!90!black} [ĐIỂM KHỞI PHÁT TỰ NHIÊN -- BẮT ĐẦU ĐỂ TRÌNH BÀY LỜI GIẢI XUÔI]}
};

\draw[arr] (b2_7) -- (b2_6);
\draw[arr] (b2_6) -- (b2_5);
\draw[arr] (b2_5) -- (b2_4);
\draw[arr] (b2_4) -- (b2_3);
\draw[arr] (b2_3) -- (b2_2);
\draw[arr] (b2_2) -- (b2_1);

\end{tikzpicture}

\end{formulabox}

\newpage

% ==============================================================
% TRANG 3: LỜI GIẢI XUÔI CHUẨN MỰC & BAREM ĐIỂM VÀO 10
% ==============================================================
\begin{solutionbox}{III. LỜI GIẢI XUÔI CHUẨN MỰC \& BAREM ĐIỂM VÀO 10}

\textbf{1. Câu a) Chứng minh tứ giác $DHEC$ nội tiếp (0,75 điểm)}
\begin{itemize}[leftmargin=5mm]
    \item Vì $AD, BE$ là các đường cao của $\triangle ABC \Rightarrow \widehat{HDC} = 90^\circ$ và $\widehat{HEC} = 90^\circ$.
    \item Tứ giác $DHEC$ có $\widehat{HDC} + \widehat{HEC} = 90^\circ + 90^\circ = 180^\circ$.
    \item Mà hai góc này ở vị trí đối nhau, nên tứ giác $DHEC$ nội tiếp đường tròn đường kính $HC$. \hfill \textit{(0,75đ)}
\end{itemize}

\vspace{0.8mm}
\textbf{2. Câu b) Chứng minh $I$ là trung điểm của $HM$ (0,75 điểm)}
\begin{itemize}[leftmargin=5mm]
    \item Vì $AM$ là đường kính của $(O)$ nên $\widehat{ABM} = 90^\circ$ và $\widehat{ACM} = 90^\circ$ (góc nội tiếp chắn nửa đường tròn).
    \item $\Rightarrow MB \perp AB$ và $MC \perp AC$.
    \item Mặt khác, $H$ là trực tâm của $\triangle ABC \Rightarrow CH \perp AB$ và $BH \perp AC$.
    \item Từ đó suy ra $BH \parallel MC$ (cùng $\perp AC$) và $CH \parallel MB$ (cùng $\perp AB$).
    \item Do đó tứ giác $BHCM$ là hình bình hành.
    \item Vì $OI \perp BC$ tại $I$ nên $I$ là trung điểm của dây cung $BC$.
    \item Trong hình bình hành $BHCM$, hai đường chéo $BC$ và $HM$ cắt nhau tại trung điểm của mỗi đường.
    \item Mà $I$ là trung điểm của $BC$ nên $I$ cũng là trung điểm của $HM$ (đpcm). \hfill \textit{(0,75đ)}
\end{itemize}

\vspace{0.8mm}
\textbf{3. Câu c) Xác định vị trí điểm $A$ để tích $DH \cdot DA$ lớn nhất (1,00 điểm)}
\begin{itemize}[leftmargin=5mm]
    \item \textbf{Bước 1: Chứng minh $DH = DK$ (0,25đ)}
    \begin{itemize}
        \item Kéo dài $AD$ cắt đường tròn $(O)$ tại điểm thứ hai $K$ ($K \neq A$).
        \item Ta có $\widehat{KBC} = \widehat{KAC}$ (hai góc nội tiếp cùng chắn $\stackrel{\frown}{KC}$).
        \item Mặt khác $\widehat{HBD} = \widehat{CAD}$ (cùng phụ với $\widehat{ACB}$).
        \item Suy ra $\widehat{KBD} = \widehat{HBD}$, do đó $BD$ vừa là đường cao vừa là phân giác trong $\triangle HBK \Rightarrow \triangle HBK$ cân tại $B$.
        \item Suy ra $BD$ là đường trung trực của $HK \Rightarrow DH = DK$.
    \end{itemize}
    \item \textbf{Bước 2: Sử dụng tam giác đồng dạng suy ra phương tích (0,25đ)}
    \begin{itemize}
        \item Khi đó: $DH \cdot DA = DK \cdot DA$.
        \item Xét $\triangle DBK$ và $\triangle DAC$ có: $\widehat{BDK} = \widehat{ADC} = 90^\circ$; $\widehat{BKD} = \widehat{ACD}$ (cùng chắn $\stackrel{\frown}{AB}$).
        \item $\Rightarrow \triangle DBK \sim \triangle DAC$ (g-g) $\Rightarrow \dfrac{DK}{DC} = \dfrac{DB}{DA} \Rightarrow DK \cdot DA = DB \cdot DC \Rightarrow DH \cdot DA = DB \cdot DC$.
    \end{itemize}
    \item \textbf{Bước 3: Áp dụng BĐT Cauchy (0,25đ)}
    \begin{itemize}
        \item Áp dụng bất đẳng thức Cauchy cho hai đoạn thẳng dương $DB$ và $DC$:
        $$DB \cdot DC \le \left(\dfrac{DB + DC}{2}\right)^2 = \left(\dfrac{BC}{2}\right)^2 = \dfrac{BC^2}{4}$$
        \item Vì dây cung $BC$ cố định nên $\dfrac{BC^2}{4}$ không đổi, do đó giá trị lớn nhất của $DH \cdot DA$ là $\dfrac{BC^2}{4}$.
    \end{itemize}
    \item \textbf{Bước 4: Xác định vị trí dấu bằng (0,25đ)}
    \begin{itemize}
        \item Dấu “=” xảy ra $\iff DB = DC \iff D \equiv I \iff \triangle ABC$ cân tại $A \iff A$ là điểm chính giữa của cung lớn $BC$. \hfill \textit{(1,00đ)}
    \end{itemize}
\end{itemize}

\end{solutionbox}

\newpage

% ==============================================================
% TRANG 4: BÌNH LUẬN SƯ PHẠM & BÀI TẬP PHÁT TRIỂN TƯƠNG ĐƯƠNG
% ==============================================================
\begin{pedagogybox}{IV. BÌNH LUẬN SƯ PHẠM: TƯ DUY KẺ ĐƯỜNG PHỤ KÉO DÀI ĐƯỜNG CAO}
\textbf{Làm thế nào để học sinh tự nghĩ ra kéo dài $AD$ cắt $(O)$ tại $K$?}
\begin{enumerate}[leftmargin=5mm]
    \item \textbf{Thầy hỏi 1:} \textit{"Đoạn thẳng $DH$ và $DA$ có đặc điểm gì về vị trí trên hình?"}\\
    $\to$ \textbf{Trò trả lời:} $A$ nằm trên $(O)$, nhưng $H$ và $D$ lại nằm lơ lửng bên trong tam giác, không thuộc đường tròn nào.
    \item \textbf{Thầy hỏi 2:} \textit{"Muốn gắn đoạn thẳng $DA$ và $DH$ vào đường tròn $(O)$ để dùng phương tích cát tuyến, ta nên làm gì?"}\\
    $\to$ \textbf{Trò trả lời:} Kéo dài $AD$ cắt $(O)$ tại $K$. Khi đó ta có ngay cát tuyến $AK$ cắt dây cung $BC$ tại $D$, lập tức cho ra $DK \cdot DA = DB \cdot DC$!
    \item \textbf{Thầy hỏi 3:} \textit{"Điểm $K$ mới sinh ra có mối quan hệ gì với trực tâm $H$?"}\\
    $\to$ \textbf{Trò trả lời:} $K$ và $H$ đối xứng nhau qua cạnh $BC$ vì $BC$ là trung trực của $HK \implies DH = DK$.
    \item \textbf{Thầy hỏi 4:} \textit{"Tích $DB \cdot DC$ có đặc điểm gì đặc biệt khi $A$ di động?"}\\
    $\to$ \textbf{Trò trả lời:} Hai đoạn $DB$ và $DC$ có tổng $DB + DC = BC$ cố định! Dùng BĐT Cauchy: Tích của hai số có tổng không đổi sẽ lớn nhất khi hai số đó bằng nhau ($DB = DC$).
\end{enumerate}
\end{pedagogybox}

\vspace{1.5mm}

\begin{exercisebox}{V. BÀI TẬP PHÁT TRIỂN TƯƠNG ĐƯƠNG (LUYỆN TẬP TẠI LỚP / VỀ NHÀ)}
\textbf{Bài 1 (Trích Đề 27 - Tuyển tập 55 Đề):} Cho đường tròn $(O)$, bán kính $R$ ($R>0$) và dây cung $BC$ cố định. Một điểm $A$ chuyển động trên cung lớn $BC$. Kẻ các đường cao $AD, BE, CF$ cắt nhau tại trực tâm $H$.
\begin{enumerate}[label=\alph*)]
    \item Chứng minh $H$ là tâm đường tròn nội tiếp của tam giác $DEF$ (tam giác Orthic).
    \item Tìm vị trí của $A$ trên $(O)$ để diện tích tam giác $HBC$ đạt giá trị lớn nhất.
\end{enumerate}

\vspace{1mm}
\textbf{Bài 2 (Trích Đề 10 - Tuyển tập 55 Đề):} Cho đường tròn $(O; R)$ và đường thẳng $d$ không cắt $(O)$. Từ điểm $M$ trên $d$ kẻ hai tiếp tuyến $MC, MD$ với $(O)$. Đường thẳng qua $O$ vuông góc với $OM$ cắt tia $MC, MD$ tại $P$ và $Q$.
\begin{enumerate}[label=\alph*)]
    \item Chứng minh $OK \cdot OM = R^2$ (với $K = OM \cap CD$).
    \item Tìm vị trí điểm $M$ trên $d$ để diện tích tam giác $MPQ$ đạt giá trị nhỏ nhất.
\end{enumerate}
\end{exercisebox}

\end{document}
